随着深度学习技术的快速发展，计算机视觉领域取得了显著进步。然而，传统的深度学习方法通常需要大量标注数据进行训练，而在细粒度分类任务中，获取充足的标注数据往往成本高昂。近年来，视觉-语言预训练模型（Vision-Language Pre-trained Models）的出现为这一问题提供了新的解决思路。

\subsection{研究动机}

细粒度图像分类旨在区分同一基础类别下的不同子类别，例如区分不同品种的猫狗。与普通分类任务相比，细粒度分类要求模型能够捕捉细微的视觉差异。传统方法通常依赖于精心设计的特征提取器或大量的标注数据，而基于预训练模型的提示学习方法则提供了一种更加高效的替代方案。

本文聚焦于猫狗品种的细粒度分类问题。猫狗作为最常见的宠物，其品种识别在宠物管理、兽医诊断等领域具有重要的实际应用价值。然而，不同品种之间的外观差异往往非常细微，这对分类算法提出了更高的要求。

\subsection{研究目标}

本研究的主要目标包括：

\begin{enumerate}
\item 设计一种自适应的提示词学习方法，能够根据输入图像动态调整提示词内容；
\item 引入可学习的难度加权机制，使模型能够自动关注难以分类的样本；
\item 在Oxford-IIIT Pets数据集上验证所提方法的有效性。
\end{enumerate}

\subsection{主要贡献}

本文的主要贡献如下：

\begin{enumerate}
\item 在 CoCoOp 图像条件提示框架基础上实现了难样本加权扩展，并在项目中以 DynamicPromptTrainer 训练器形式组织；
\item 设计了可学习的难度权重计算器，使模型能够自动学习对不同难度样本的加权策略；
\item 在Oxford-IIIT Pets数据集上的实验表明，本文方法在多数few-shot设置下优于CoOp和CoCoOp基线方法。
\end{enumerate}
